
\section{Probabilistic Boolean Circuits}\label{sec:probboolcircuits}
Probabilistic Boolean circuits were introduced in \cite{piedeleu2025boolean} as a string-diagrammatic language for finite probabilistic programming. In this section, we recall their syntax and their semantics. In the following section, we provide a complete axiomatization by encoding them into tape diagrams.

Our presentation is staged. We begin with Boolean circuits (Section~\ref{ssec:boolean}), then move to partial Boolean circuits (Section~\ref{ssec:partialboolean}), and finally illustrate probabilistic Boolean circuits (Section~\ref{ssec:probabilisticboolean}). We also provide a complete axiomatization for partial Boolean circuits, which, to the best of our knowledge, has not previously appeared in the literature.



\subsection{Boolean Circuits}\label{ssec:boolean}
Consider the monoidal signature with a single sort, $\sort = \{ A\}$, and  generators 
\[ \SigB \defeq \{\Andgate  , \Notgate , \Flip{1},   \CBcopier, \CBdischarger \} \text{.}\]  
Arities and coarities are determined by the number of ports on the left and on the right: for instance, $\Andgate$ has arity $A^2=AA$ and coarity $A$, while $\CBdischarger $ has arity $A$ and coarity $A^0 =1$. 

The arrows of $\Diag{\SigB}$, namely the category of string diagrams generated by the monoidal signature $\SigB$ (see Section~\ref{sec:tapediagrams}), are interpreted as Boolean functions. 
The first three generators of $\SigB$ correspond to the operations and constants of Boolean algebras, namely $\wedge$, $\neg$, and $1$. 
The generator \emph{copy}, denoted by $\CBcopier$, takes a Boolean signal as input and produces two identical outputs. 
The generator \emph{discard}, denoted by $\CBdischarger$, takes a Boolean signal as input and discards it.


Formally, the semantics of these diagrams is given by the monoidal functor  $\osemB{-}\colon \Diag{{\SigB}} \to \Sets$ defined on objects as $A^n \mapsto 2^n$ where $2$ is the set of Booleans  $\{0,1\}$; for arrows it is defined by first interpreting the generators as functions in $\Sets$
  \[
    \begin{array}{rclrclrcl}
      \osemB{\Andgate}\colon 2 \times 2& \to & 2  & \osemB{\Notgate}\colon  2& \to & 2 & \osemB{\Flip{1}} \colon 1 & \to & 2 \\
      (x,y)&\mapsto& x\wedge y &  x & \mapsto & \neg x & \bullet &\mapsto&1\\[4pt]
       \osemB{\CBcopier} \colon 2  & \to & 2\times 2 & \osemB{\CBdischarger} \colon 2 & \to & 1 & %\osemB{\Flip{p}} \colon 1 & \to & 2  
       \\
      x &\mapsto&(x,x) & x &\mapsto&{\bullet} & %\bullet &\mapsto& \delta_1 +_p \delta_0
    \end{array}
  \]
and then inductively by means of the structure of the monoidal category $(\Sets, \otimes ,\uno)$: for instance, $\osemB{c_1\otimes c_2}=\osemB{c_1}\otimes \osemB{c_2}$ where the second $\otimes$ is the standard cartesian product of functions. 

The monoidal theory $\ThB$ consists of the signature $\SigB$ together with the axioms listed in Figure~\ref{tab:booleanalgebra}. 
The equalities in the first block are the standard axioms of Boolean algebras. 
The constant $\Flip{0}$ and the $\Orgate$ gate are defined as expected (see Figure~\ref{table:multiplexer}). 
Note that, in the string-diagrammatic setting, the structural maps $\CBcopier$ and $\CBdischarger$ must be made explicit.
The equalities in the second block state that $\CBcopier$ and $\CBdischarger$ form a commutative comonoid. 
The equalities in the third and fourth blocks impose naturality of $\CBcopier$ and $\CBdischarger$, respectively.

%The monoidal theory $\ThB$ consists of the signature $B$ and the axioms in Figure~\ref{tab:booleanalgebra}. The equalities on the first block are the standard axioms of Boolean algebras, where the constant $\Flip{0}$ and the or gate $\Orgate$ are defined as expected in Figure~\ref{table:multiplexer}: note that the use of string diagrams forces to make 
%explicit use of $\CBcopier$ and $\CBdischarger$; these form a comonoid by the equalities in the second block; the equalities in the third and fourth blocks impose naturality of $\CBcopier$ and $\CBdischarger$, respectively. 

The congruence with respect to $;$ and $\otimes$ generated by the axioms in Figure~\ref{tab:booleanalgebra} is denoted by $\eqB$. 
We write $\Diag{\ThB}$ for the quotient of $\Diag{\SigB}$ by $\eqB$.

As expected, the equalities in Figure~\ref{tab:booleanalgebra} provide a sound and complete axiomatization of the equality induced by $\osemB{-}$.

%The functor $\osem{-}$ trivially extends to a monoidal functor $\Diag{\ThB} \to \KlD$, since the axioms in Figure~\ref{tab:booleanalgebra} are sound with respect to the semantics.
%For Boolean circuits, we recall the following completeness result from \cite[Theorem 3.2]{piedeleu2025boolean}.
\begin{theorem}\label{thm:completenessbooleancircuits}
For all  $c,d\in \Diag{\SigB}$, $\osemB{c}=\osemB{d}$ iff $c\eqB d$.
\end{theorem}
\begin{proof}
See e.g. \cite[Theorem 3.2]{piedeleu2025boolean}.
\end{proof}



\begin{figure}
  \scalebox{0.8}{%
  \begin{tabular}{ccccc}

  % =====================
  % ---- B axioms ----
  % =====================

  \mylabelbis{ax:B1}{B1}%
  $\tikzfig{tapes/cipriano/axB1left}\raisebox{-1.5ex}{\,$\stackrel{\eqref{ax:B1}}{=}$\,}\tikzfig{tapes/cipriano/axB1right}$
  & &
  \mylabelbis{ax:B2l}{B2l}\mylabelbis{ax:B2r}{B2r}%
  $\tikzfig{tapes/cipriano/axB2left}\raisebox{-1.5ex}{\,$\stackrel{\eqref{ax:B2l}}{=}$\,}\tikzfig{tapes/cipriano/axB2middle}\raisebox{-1.5ex}{\,$\stackrel{\eqref{ax:B2r}}{=}$\,}\tikzfig{tapes/cipriano/axB2right}$
  & &
  \mylabelbis{ax:B3}{B3}%
  $\tikzfig{tapes/cipriano/axB3left}\raisebox{-1.5ex}{\,$\stackrel{\eqref{ax:B3}}{=}$\,}\tikzfig{tapes/cipriano/axB3right}$
  \\

  \mylabelbis{ax:B4}{B4}%
  $\tikzfig{tapes/cipriano/axB4left}\raisebox{-1.5ex}{\,$\stackrel{\eqref{ax:B4}}{=}$\,}\tikzfig{tapes/cipriano/axB4right}$
  & &
  \mylabelbis{ax:B5}{B5}%
  $\tikzfig{tapes/cipriano/axB5left}\raisebox{-1.5ex}{\,$\stackrel{\eqref{ax:B5}}{=}$\,}\tikzfig{tapes/cipriano/axB5right}$
  & &
  \mylabelbis{ax:B6}{B6}%
  $\tikzfig{tapes/cipriano/axB6left}\raisebox{-1.5ex}{\,$\stackrel{\eqref{ax:B6}}{=}$\,}\tikzfig{tapes/cipriano/axB6right}$
  \\

  & &
  \mylabelbis{ax:B7}{B7}%
  $\tikzfig{tapes/cipriano/axB7left}\raisebox{-1.5ex}{\,$\stackrel{\eqref{ax:B7}}{=}$\,}\tikzfig{tapes/cipriano/axB7right}$
  & &
  \\

  \midrule

  % =====================
  % ---- A axioms ----
  % =====================

  \mylabelbis{ax:A1}{B8}%
  $\tikzfig{tapes/cipriano/axA1left}\raisebox{-1.5ex}{\,$\stackrel{\eqref{ax:A1}}{=}$\,}\tikzfig{tapes/cipriano/axA1right}$
  & &
  \mylabelbis{ax:A2l}{B9l}\mylabelbis{ax:A2r}{B9r}%
  $\tikzfig{tapes/cipriano/axA2left}\raisebox{-1.5ex}{\,$\stackrel{\eqref{ax:A2l}}{=}$\,}\tikzfig{tapes/cipriano/axA2middle}\raisebox{-1.5ex}{\,$\stackrel{\eqref{ax:A2r}}{=}$\,}\tikzfig{tapes/cipriano/axA2right}$
  & &
  \mylabelbis{ax:A3}{B10}%
  $\tikzfig{tapes/cipriano/axA3left}\raisebox{-1.5ex}{\,$\stackrel{\eqref{ax:A3}}{=}$\,}\tikzfig{tapes/cipriano/axA3right}$
  \\

  \midrule

  % =====================
  % ---- C axioms ----
  % =====================

  \mylabelbis{ax:C1}{B11}%
  $\tikzfig{tapes/cipriano/axC1left}\raisebox{-0.6ex}{\,$\stackrel{\eqref{ax:C1}}{=}$\,}\tikzfig{tapes/cipriano/axC1right}$
  & &
  \mylabelbis{ax:C2}{B12}%
  $\tikzfig{tapes/cipriano/axC2left}\raisebox{-0.6ex}{\,$\stackrel{\eqref{ax:C2}}{=}$\,}\tikzfig{tapes/cipriano/axC2right}$
  & &
  \mylabelbis{ax:C3}{B13}%
  $\tikzfig{tapes/cipriano/axC3left}\raisebox{-0.6ex}{\,$\stackrel{\eqref{ax:C3}}{=}$\,}\tikzfig{tapes/cipriano/axC3right}$
  \\

  \midrule

  % =====================
  % ---- D axioms ----
  % =====================

  \mylabelbis{ax:D1}{B14}%
  $\tikzfig{tapes/cipriano/axD1left}\raisebox{-0.6ex}{\,$\stackrel{\eqref{ax:D1}}{=}$\,}\tikzfig{tapes/cipriano/axD1right}$
  & &
  \mylabelbis{ax:D2}{B15}%
  $\tikzfig{tapes/cipriano/axD2left}\raisebox{-0.6ex}{$\stackrel{\eqref{ax:D2}}{=}$}\tikzfig{tapes/cipriano/axD2right}$
  & &
  \mylabelbis{ax:D3}{B16}%
  $\tikzfig{tapes/cipriano/axD3left}\raisebox{-0.6ex}{$\stackrel{\eqref{ax:D3}}{=}$}
  \raisebox{-0.5em}{\tikzfig{empty}}$
  \\

  \end{tabular}}%
  \caption{The monoidal theory of Boolean algebras.}
  \label{tab:booleanalgebra}
\end{figure}
Moreover, one can characterise exactly the image of $\Diag{\SigB}$ through $\osemB{-}$. Let $\Sets_2$ be the full subcategory of sets where objects are powers of the set $2$, i.e.,  $2^n$ for all $n \in \mathbb{N}$.


\begin{proposition}\label{prop:foolboolean}
The monoidal functor  $\osemB{-}\colon \Diag{{\SigB}} \to \Sets$ factors as
\[\xymatrix{ \Diag{\SigB} \ar@{->>}[r]& \Diag{\ThB} \ar[r]^{\cong}& \Sets_2 \ar@{^{(}->}[r]& \Sets }\]
where the leftmost and rightmost functor are the obvious quotients and injections and the central arrow is a monoidal isomorphism.
\end{proposition}
\begin{proof}
Faithfulness of the central functor is a direct consequence of Theorem~\ref{thm:completenessbooleancircuits}. Fullness follows from the fact that every function $f\colon 2^n \to 2^m$ corresponds to a vector of $m$ Boolean formulas in $n$ variables, which can be easily encoded as a string diagram in $\Diag{\ThB}$ using Boolean operators, $\CBcopier$ and $\CBdischarger$.
\end{proof}



We conclude this section by fixing some syntactic sugar, collected in Figure~\ref{table:multiplexer}. 
The \emph{xnor} gate $\xnor$  takes two Boolean inputs and outputs $1$ if they are equal and $0$ otherwise. 
The \emph{multiplexer} $\tikzfig{tapes/cipriano/multi}$ takes three inputs and returns one output: if the first input is $1$, then it outputs the second input, otherwise the third input. 
The $m$-ary multiplexer $\Ifgatem$ works similarly, but it takes $2m+1$ inputs: if the first input is $1$, then it outputs the first $m$ inputs, otherwise it outputs the second $m$ inputs. 
Figure~\ref{table:multiplexer} also defines $m$-ary discard $\CBdischargerm$ and copier $\mcopier$. Note that, for the sake of readability of string diagrams, we label wires by  some natural number $n$ as an abbreviation for $A^n$.

\begin{comment}
\begin{lemma}\label{lemma:multiplexer}\marginpar{Puoi fare questi come string diagrams? Non cancellare i vecchi che forse ci servono se facciamo la versione conferenza}
  Let $c$ be a diagram in $\Diag{\SigB}[n,m]$. The following derived laws hold:
  \begin{enumerate}[label={}, leftmargin=*]

    \item[(D1)]\mylabelbis{ax:M1}{D1}
      $(\Flip{0}\otimes \id{m}\otimes \id{m});\Ifgatem
      \eqB \CBdischargerm\otimes \id{m}$;

    \item[(D2)]\mylabelbis{ax:M2}{D2}
      $(\Flip{1}\otimes \id{m}\otimes \id{m});\Ifgatem
      \eqB \id{m}\otimes \CBdischargerm$;

    \item[(D3)]\mylabelbis{ax:M3}{D3}
      $(\id{}\otimes \id{} \otimes \Flip{0});\Ifgate
      \eqB \Andgate$;

    \item[(D4)]\mylabelbis{ax:N1}{D4}
      $c;\CBdischargerm \eqB \CBdischargern$;

    \item[(D5)]\mylabelbis{ax:N2}{D5}
      $c;\mcopier \eqB \ncopier;(c\otimes c)$.

  \end{enumerate}
\end{lemma}
\end{comment}


\begin{lemma}\label{lemma:multiplexer}
 Let $c$ be a diagram in $\Diag{\SigB}[n,m]$. The following derived laws hold:

\mylabelbis{ax:M1}{D1}
\mylabelbis{ax:M2}{D2}
\mylabelbis{ax:M3}{D3}
\mylabelbis{ax:N1}{D4}
\mylabelbis{ax:N2}{D5}
{\setlength{\abovedisplayskip}{0.3em}
 \setlength{\belowdisplayskip}{0.3em}
\[
\begin{array}{c@{\qquad}c@{\qquad}c}
%
\raisebox{-0.3em}{\scalebox{0.8}{\tikzfig{tapes/cipriano/D1left}}} 
\stackrel{\eqref{ax:M1}}{\eqB}
\tikzfig{tapes/cipriano/D1right}
%
&
\raisebox{-0.3em}{\scalebox{0.8}{\tikzfig{tapes/cipriano/D2left}}}
\stackrel{\eqref{ax:M2}}{\eqB}
\tikzfig{tapes/cipriano/D2right}
%
&
\raisebox{-0.3em}{\scalebox{0.8}{\tikzfig{tapes/cipriano/D3left}}}
\stackrel{\eqref{ax:M3}}{\eqB}
\tikzfig{tapes/cipriano/D3right}
%
\\[0.5em]
%
\tikzfig{tapes/cipriano/D4left}
\stackrel{\eqref{ax:N1}}{\eqB}
\tikzfig{tapes/cipriano/D4right}
%
&
\tikzfig{tapes/cipriano/D5left}
\stackrel{\eqref{ax:N2}}{\eqB}
\tikzfig{tapes/cipriano/D5right}
%
&
%
\end{array}
\]}
\end{lemma}




\begin{figure}[t]
\centering
\resizebox{0.96\linewidth}{!}{%
\begin{tabular}{@{}l@{}}

% ---- Row 1 ----
$\Orgate \defeq (\Notgate \otimes \Notgate); \Andgate; \Notgate 
\qquad 
\Flip{0} \defeq \Flip{1}; \Notgate
\qquad
\tikzfig{tapes/cipriano/multi} 
\defeq 
\tikzfig{tapes/cipriano/multiexpl}$ \\[4pt]

% ---- Row 2 ----
$\tikzfig{tapes/cipriano/xnorgate}
\qquad
\raisebox{-1em}{\tikzfig{tapes/cipriano/ncopierzero}} 
\raisebox{-0.5em}{$\defeq$}
\raisebox{-0.5em}{\tikzfig{empty}}
\qquad 
\raisebox{-1em}{\tikzfig{tapes/cipriano/ncopier}}$ \\[6pt]

% ---- Row 3 ----
$\raisebox{-1.2em}{\tikzfig{tapes/cipriano/mmultiplexerleftzero}} 
\defeq 
\tikzfig{tapes/cipriano/mmultiplexerrightzero}
\qquad 
\raisebox{-1.2em}{\tikzfig{tapes/cipriano/mmultiplexerleft}} 
\defeq 
\tikzfig{tapes/cipriano/mmultiplexerright}$ \\[6pt]

% ---- Row 4 ----
$\mFlip{1}{0} \defeq \tikzfig{empty}
\qquad
\mFlip{1}{m+1} \defeq \Flip{1} \otimes \mFlip{1}{m}
\qquad
\mCBdischarger{0} \defeq \tikzfig{empty}
\qquad
\mCBdischarger{m+1} \defeq \CBdischarger\otimes \mCBdischarger{m}$

\end{tabular}
}
\caption{Definitions of OR gate, $\Flip{0}$ and multiplexer (top); definitions of XNOR gate and inductive definition of $n$-ary copier (middle); inductive definitions of $m$-ary multiplexer (third row); definitions of $\Flipzero{1}$, $\Flipm{1}$ and $n$-ary dischargers (bottom).}
\label{table:multiplexer}
\end{figure}



%
%
% SUBSECTION
%
%
\subsection{Partial Boolean Circuits}\label{ssec:partialboolean}
Now consider the monoidal signature obtained by adding to $\SigB$ the generator \emph{cocopy} $\CBcocopier$ of arity $A^2$ and coarity $A$. 
\[ PB\defeq \SigB \cup\{\CBcocopier\}\text{}\] 
Intuitively, $\CBcocopier$ is the dual of $\CBcopier$: it compares two inputs and, if these are equal, it outputs that value; otherwise, it produces no output. Thus, string diagrams in $\Diag{PB}$ represent \emph{partial} Boolean functions. 

To formally define their semantics, let us first fix $(\Cat{Par}, \otimes, 1)$ to be the symmetric monoidal category of sets and partial functions, where the monoidal product $\otimes$ is the cartesian product of sets and, for all partial functions $f_1\colon X_1 \to Y_1$ and $f_2\colon X_2 \to Y_2$, $f_1\otimes f_2 \colon X_1 \times X_2 \to Y_1 \times Y_2$ is defined as
\begin{equation}\label{eq:productpar}
f_1\otimes f_2 (x_1,x_2) \defeq \begin{cases} (\,f_1(x_1), \, f_2(x_2)\,) & \text{ if }f_1(x_1)\neq \bot\text{, } f_2(x_2)\neq \bot \\
 \bot & \text{ otherwise}\end{cases}
\end{equation}
for all $(x_1,x_2)\in X_1\times X_2$. Above, and in the rest of the paper, we write $f(x)=\bot$ to say that a partial function $f\colon X \to Y$ is undefined for some $x\in X$. 
One can readily check that the obvious injection $\Pa\colon \Sets \to \Cat{Par}$ is a symmetric monoidal functor.


Now, the semantics of diagrams in $\Diag{PB}$ is given by the monoidal functor $\osemPB{-}\colon \Diag{PB} \to \Cat{Par}$ defined on objects as $A^n \mapsto 2^n$, like $\osemB{-}$; 
for arrows it is defined by first interpreting generators $c \in \SigB$ as $\osemPB{c}\defeq \Pa(\osemB{c})$, the generator $\CBcocopier $ as
  \begin{equation}\label{eq:semcocopier}
    \begin{array}{rcl}
      \osemPB{\CBcocopier}\colon 2 \times 2& \to & 2   \\
      (x,y)&\mapsto& \begin{cases} x &\text{if } x=y\\ \bot & \text{else}\end{cases}    \end{array}
  \end{equation}
and then inductively by means of the structure of the monoidal category $(\Cat{Par}, \otimes ,\uno)$: for instance $\osemPB{c_1\otimes c_2}=\osemPB{c_1}\otimes \osemPB{c_2}$ where the second $\otimes$ is the one in \eqref{eq:productpar}. %A simple inductive argument confirms that $\osemPB{-}$ is an extension of $\osemB{-}$, i.e., the following diagram commutes
%\[
%\xymatrix{\Diag{PB} \ar[r]^{\osemPB{-}} & \Cat{Par}\\ 
%\Diag{B} \ar[r]_{\osemB{-}} \ar@{>->}[u] & \Sets \ar@{>->}[u]_P
%}
%\]
For example, the semantics of the following diagrams
\begin{equation}\label{def:bottom}
  \raisebox{0cm}{\tikzfig{tapes/cipriano/defcoflip}}
\end{equation}
is illustrated below.
  \[
    \begin{array}{rclrclrcl}
      \osemPB{\coflip{1} }\colon 2 & \to & 1  & \osemPB{\coflip{0} }\colon 2 & \to & 1 & \osemPB{\Flip{\bot}}\colon 1 & \to & 2 \\
      x&\mapsto& \begin{cases} \bullet &\text{if } x=1\\ \bot & \text{else}\end{cases}    
      &
      x&\mapsto& \begin{cases} \bullet &\text{if } x=0\\ \bot & \text{else}\end{cases}    
      &
      \bullet &\mapsto & \bot     
      \end{array}
  \]






In this section, we present a complete axiomatization of the equality induced by $\osemPB{-}$, which, to the best of our knowledge, has not been previously established.
The axioms are illustrated in Figure~\ref{tab:partialbooleanalgebra}: the equalities in the first row assert that $\CBcocopier$ is associative, commutative and idempotent; in the second row, there are the usual Frobenius equalities ruling the interaction of $\CBcocopier$ with $\CBcopier$; the two axioms in the third row are more peculiar: the rightmost states that $x\wedge y = 1$ iff $x=y=1$; the leftmost provides a decomposition of  $\CBcocopier$ in terms of $\Andgate$, $\coflip{1}$ and the xnor gate \tikzfig{tapes/cipriano/xnorgatealone}.



%%%%%%%NUOVA TABELLA CON TAG%%%%%%%%%%
\begin{figure}
\centering
\scalebox{0.8}{
{
\setlength{\tabcolsep}{12pt}      % spazio orizzontale
\setlength{\extrarowheight}{9pt}  % spazio verticale aggiuntivo per OGNI riga
\renewcommand{\arraystretch}{1.0} % qui può anche restare 1, extrarowheight fa il grosso

\begin{tabular}{ccc}

% Prima riga
\mylabelbis{ax:F1}{P1}%
$\tikzfig{tapes/cipriano/axF1left}
 \raisebox{-0.6em}{\,\,$\stackrel{\eqref{ax:F1}}{=}$\,\,}
 \tikzfig{tapes/cipriano/axF1right}$
&
\mylabelbis{ax:F2}{P2}%

$\tikzfig{tapes/cipriano/axF2left}
 \raisebox{-0.6em}{\,$\stackrel{\eqref{ax:F2}}{=}$\,}
 \tikzfig{tapes/cipriano/axF2right}$
&
\mylabelbis{ax:F3}{P3}%
$\tikzfig{tapes/cipriano/axF3left}
 \raisebox{-0.6em}{\,\,$\stackrel{\eqref{ax:F3}}{=}$\,\,}
 \tikzfig{tapes/cipriano/axF3right}$
\\[1.0em] % spazio extra tra riga 1 e riga 2

% Seconda riga (F4 e F5)
\multicolumn{3}{c}{%
  \mylabelbis{ax:F4}{P4}%
  \mylabelbis{ax:F5}{P5}%
  $\tikzfig{tapes/cipriano/axF4left}
    \,\stackrel{\eqref{ax:F4}}{=}\,
    \tikzfig{tapes/cipriano/axF4right}
    \,\stackrel{\eqref{ax:F5}}{=}\,
    \tikzfig{tapes/cipriano/axF5}$%
} \\[1.0em] % spazio extra tra riga 2 e riga 3

% Terza riga
\multicolumn{3}{c}{%
  \mylabelbis{ax:F6}{P6}%
  $\tikzfig{tapes/cipriano/axF6left}
   \, \stackrel{\eqref{ax:F6}}{=}\,
    \tikzfig{tapes/cipriano/axF6right}$%
  \hspace{4em}%
  \mylabelbis{ax:F7}{P7}%
  $\tikzfig{tapes/cipriano/axF7left}
    \,\stackrel{\eqref{ax:F7}}{=}\,
    \tikzfig{tapes/cipriano/axF7right}$%
}

\end{tabular}
}
}
\caption{The monoidal theory of partial Boolean algebras.}
\label{tab:partialbooleanalgebra}
\end{figure}


We write $\ThPB$ for the monoidal theory consisting of the signature $PB$ and the axioms in Figures~\ref{tab:booleanalgebra} and~\ref{tab:partialbooleanalgebra};  $\eqPB$ for the generated congruence; $\Diag{\ThPB}$ for the quotient of $\Diag{PB}$ by $\eqPB$ and $Q_\ThPB\colon \Diag{ \SigPB} \to \Diag{\ThPB}$ the functor mapping each diagram into its $\eqPB$ equivalence class. Simple computations confirm that:

\begin{proposition}[Soundness]\label{prop:soundnesspartialb}
For all $c,d\in \Diag{PB}$, if $c\eqPB d$, then $\osemPB{c}=\osemPB{d}$.
\end{proposition}




%%%%%%FRASE GIA MESSA NEL PARAGRAFO PRECEDENTE%%%%%%%%%%%
%The equalities on the first block are the standard axioms of Boolean algebras, where the constant $\Flip{0}$ and the or gate $\Orgate$ are defined as expected in Table \ref{tab:partialbooleanalgebra}: note that the use of string diagrams forces to make explicit use of $\CBcopier$ and $\CBdischarger$; these form a comonoid by the equalities in the second block; the equalities in the third and fourth blocks impose naturality of $\CBcopier$ and $\CBdischarger$, respectively. 






%The functor $\osemPB{-}$ trivially extends to a monoidal functor $\Diag{\ThPB} \to \Cat{Par}$, since the axioms in Figure~\ref{tab:booleanalgebra} and \ref{tab:partialbooleanalgebra} are sound with respect to the semantics.



% We refer to the arrows of the category  $\Diag{\ThPB}$, obtained as the quotient of $\Diag{PB}$  by the axioms in Figure~\ref{tab:booleanalgebra} and \ref{tab:partialbooleanalgebra}, 
% as \textit{partial boolean circuits}.




%%%%%%%TOLTO PERCHè GIà DEFINITI PIù SOPRA CON ANCHE INTERPETAZIONE SEMANTICA %%%%%%%%%%
%To achieve this, we first introduce some syntactic sugar and prove some auxiliary properties.
%First consider the following circuits:
%\begin{center}
%  \tikzfig{tapes/cipriano/defcoflip}
%\end{center}

The remainder of this section is devoted to proving the converse implication, namely completeness. 
Although this fact is not required for our proof, it is worth noting that the axiom \textsc{F8} in~\cite{piedeleu2025boolean}--corresponding to \eqref{ax:F8} in the following lemma--can be derived within $\ThPB$. 

\begin{lemma}\label{lemma:failureequalities}
  The following derived laws hold:
  \begin{center}
    \mylabelbis{ax:F9}{D6}%
    $\tikzfig{tapes/cipriano/axF9left}
      \stackrel{\eqref{ax:F9}}{\eqPB}
      \tikzfig{tapes/cipriano/axF9right}$
    \qquad
    \mylabelbis{ax:F8}{D7}%
    $\tikzfig{tapes/cipriano/axF8left}
      \stackrel{\eqref{ax:F8}}{\eqPB}
      \tikzfig{tapes/cipriano/axF8right}$
    \qquad
    \mylabelbis{ax:F10l}{D8l}%
    \mylabelbis{ax:F10r}{D8r}%
    $\tikzfig{tapes/cipriano/axF10left}
      \stackrel{\eqref{ax:F10l}}{\eqPB}
      \tikzfig{tapes/cipriano/axF10middle}
      \stackrel{\eqref{ax:F10r}}{\eqPB}
      \tikzfig{tapes/cipriano/axF10right}$
  \end{center}
\end{lemma}
\begin{proof}
 \eqref{ax:F9} follows from the equalities  
 \begin{center} \tikzfig{tapes/cipriano/lemmafail1}
  \end{center}
  \begin{center}
  \tikzfig{tapes/cipriano/lemmafail2}
  \end{center}
   \eqref{ax:F8} follows from the equalities
  \begin{center}
  \tikzfig{tapes/cipriano/lemmafail3}
  \end{center}
  \begin{center}
  \tikzfig{tapes/cipriano/lemmafail4}
  \end{center}
  \begin{center}
  \tikzfig{tapes/cipriano/lemmafail5}
  \end{center}
  \begin{center}
  \tikzfig{tapes/cipriano/lemmafail5mezzo}
  \end{center}
   \eqref{ax:F10l} follows from the equalities
  \begin{center}
  \tikzfig{tapes/cipriano/lemmafail6}
  \end{center}
  \eqref{ax:F10r} can be proved with a similar argument.
\end{proof}




\begin{comment}
  \begin{corollary}\label{cor:failureproperty} For all $f,g\in \Diag{PB}[A^n,A^m]$, the following equality holds.\marginpar{SE NON SI USA LEVIAMOLO}
  \begin{center}
  \tikzfig{tapes/cipriano/lemmafailure}
  \end{center}
\end{corollary}
\begin{proof}
  It follows by structural induction as in \cite[Lemma 4.4]{piedeleu2025boolean}, observing that the above lemma allows to derive axiom \eqref{ax:F8} of \cite{piedeleu2025boolean}.
\end{proof}
\end{comment}

The strategy for proving completeness is as follows. 
First, we show that any diagram $c$ in $\Diag{PB}$ can be decomposed into two diagrams $D_c$ and $T_c$ in $\Diag{\SigB}$ (Proposition~\ref{prop:decompositionpartialcircuits}). 
We then appeal to the completeness result for Boolean circuits (Theorem~\ref{thm:completenessbooleancircuits}).

Intuitively, $D_c$ and $T_c$ represent, respectively, the domain and the total component of $c$. 
The diagram $D_c$ returns $1$ if the input of $c$ lies in the domain of definition of $c$, and $0$ otherwise. 
The diagram $T_c$ returns the output of $c$ when the input lies in its domain, and a vector of $0$ otherwise.

\begin{definition}
For all $c\in \Diag{PB}[A^n,A^m]$, $D_c\in \Diag{\SigB}[A^n,A]$ and $T_c\in\Diag{\SigB}[A^n,A^m]$ are inductively defined as
\begin{equation}\label{eq:defTD}  \begin{array}{rcl@{\hspace{1.5cm}}rcl}
    D_c &\defeq& \CBdischargern\, \Flip{1} &       T_c &\defeq& c \qquad \text{for all $c\in \SigB \cup \{\id{1},\id{A},\symmt{A}{A}\}$;}\\
    D_{\scalebox{0.6}{$\CBcocopier$}} &\defeq& \tikzfig{tapes/cipriano/xnorgatealone} &     T_{\scalebox{0.6}{$\CBcocopier$}} &\defeq& \Andgate\\
    D_{c;d} &\defeq &\quad \tikzfig{tapes/cipriano/Dfg} &  T_{c;d} &\defeq &\quad \tikzfig{tapes/cipriano/Tfg}\\
    D_{c\otimes d} &\defeq & \quad \tikzfig{tapes/cipriano/Dfperg} & T_{c\otimes d} &\defeq& \quad \tikzfig{tapes/cipriano/Tfperg}\\
  \end{array}
  \end{equation}
\end{definition}
Observe that the domain of $\CBcocopier$ is the $\mathrm{xnor}$ gate: it returns $1$ if and only if the two inputs are equal. 
Its total part is given by the $\wedge$ gate: it returns $1$ precisely when both inputs are $1$.
Similarly, the domain of $c \otimes d$ is given by the conjunction of the domains of $c$ and $d$. 
Its total part returns the total parts of $c$ and $d$ when the domain evaluates to $1$, and $0$ otherwise.

\begin{comment}
  \begin{lemma}\label{lemma:DTboolean}\marginpar{Non sono sicuro che questo lemma serva}
  For any partial Boolean circuit $f\colon n\to m$, it holds that $D_f$ and $T_f$ are purely Boolean, i.e., they are generated by the generators in $\SigB$ and axioms in Figure~\ref{tab:booleanalgebra}.
\end{lemma}
\begin{proof}
  We proceed by induction on the structure of $f$. For the base case, if $f$ is a Boolean gate, then the statement is trivial. If $f$ is $\CBcocopier$, then the statement follows from the definition of $D_{\scalebox{0.6}{$\CBcocopier$}}$ and $T_{\CBcocopier}$. Now, if $f$ is of the form $g;h$, then assuming that the statement holds for $g$ and $h$, i.e. $D_g,D_h,T_g,T_h$ are Boolean, we have that $D_f$ and $T_f$ are purely Boolean since they are defined by composition and tensor of Boolean circuits with the multiplexer, which is Boolean by definition. Finally, if $f$ is of the form $g\otimes h$, then assuming that the statement holds for $g$ and $h$, i.e. $D_g,D_h,T_g,T_h$ are Boolean, we have that $D_f$ and $T_f$ are purely Boolean since they are defined by composition and tensor of Boolean circuits with the AND gate and the multiplexer, which are Boolean by definition.
\end{proof}
\end{comment}




%We now prove that every partial Boolean circuit $f$ can be decomposed through its domain and total part as follows.

\begin{proposition}\label{prop:decompositionpartialcircuits}
  For all $c\in \Diag{BP}$, it holds that
  \begin{center}
  \tikzfig{tapes/cipriano/decompositionpartial}
  \end{center}
\end{proposition}
\begin{proof}
  We proceed by induction on the structure of $c$. For the base case, if $c\in \SigB \cup \{\id{1},\id{A},\symmt{A}{A}\}$, then the statement follows by
  \begin{center}
    \tikzfig{tapes/cipriano/propdecomposition1}
  \end{center} 
  If $c$ is $\CBcocopier$, then the statement is exactly axiom \eqref{ax:F6} in Figure~\ref{tab:partialbooleanalgebra}. Now, in the case $c;d$, then assuming that the statement holds for $c$ and $d$, we have
  \begin{center}
    \tikzfig{tapes/cipriano/propdecomposition2}
    \end{center}
    \begin{center}
    \tikzfig{tapes/cipriano/propdecomposition3}
    \end{center}
    \begin{center}
    \tikzfig{tapes/cipriano/propdecomposition4}
    \end{center}
    \begin{center}
    \tikzfig{tapes/cipriano/propdecomposition5}
    \end{center}
    where now the statement follows by inductive hypothesis. Finally, in the case $c\otimes d$, then assuming that the statement holds for $c$ and $d$, we have %\marginpar{Rifarei le etichette con $n',n'',m',m''$ poiche' $n$ e $m$ sono fissati nello statemnet del lemma e non sono le stesse di questa derivazione-C:Fatto}
    \begin{center}
    \tikzfig{tapes/cipriano/propdecomposition6}
    \end{center}
    \begin{center}
    \tikzfig{tapes/cipriano/propdecomposition7}
    \end{center}
    \begin{center}
    \tikzfig{tapes/cipriano/propdecomposition8}
    \end{center}
    now applying the inductive hypothesis for $c$ and $d$, we obtain the statement.
\end{proof}

%Hereafter, we write ${B}[1,A^n]$ for the set $n$-ary Boolean vectors, i.e., all those Boolean circuits in ${\Diag{PB}}[1,A^n]$ of the form $\bigotimes_{i=1}^n b_i$ for $b_i \in \{\Flip{0}, \Flip{1}\}$. \Diag{\SigB}
Before proceeding, recall that Theorem~\ref{thm:completenessbooleancircuits} implies that every $b\in \Diag{\SigB}[1,A^n]$ is $\eqB$-equal (hence $\eqPB$-equal) to a Boolean vector, i.e., a circuit of the form $\bigotimes_{i=1}^n b_i$ for $b_i \in \{\Flip{0}, \Flip{1}\}$ if $n>0$, and $\tikzfig{empty}$ if $n=0$. 

\begin{lemma}\label{lemma:complete1}
  Let $b\in \Diag{\SigB}[1,A^n]$ and $c\colon n\to m$ be a partial Boolean circuit. If $b;D_c \eqPB \Flip{0}$, then $b;T_c\eqPB\Flipm{0}$.
\end{lemma}
\begin{proof}
  We proceed by induction on the structure of $c$. For the base case, if $c\in \SigB \cup \{\id{1},\id{A},\symmt{A}{A}\}$, then $b;D_c=b;\CBdischargern\, \Flip{1}=\Flip{1}$, hence the statement is trivial. If $c$ is $\CBcocopier$, then, since $D_{\scalebox{0.6}{$\CBcocopier$}}$ is the xnor gate, $b$ must be of the form $\Flip{1}\otimes \Flip{0}$ or $\Flip{0}\otimes \Flip{1}$, hence $b;T_{\scalebox{0.6}{$\CBcocopier$}}$ is $\Flip{0}$, by axioms \eqref{ax:B2l} and \eqref{ax:B2r} in Figure~\ref{tab:booleanalgebra}. Now, in the case $c;d$, we have
  \begin{center}
    \tikzfig{tapes/cipriano/lemma1fg}
  \end{center}
  \begin{center}
    \tikzfig{tapes/cipriano/lemma1fgpt2}
  \end{center}
   Finally, in the case $c\otimes d$, since  $b\eqPB b_1\otimes b_2$ for some $b_1\in \Diag{\SigB}[1,A^{n'}]$ and $b_2\in \Diag{\SigB}[1,A^{n''}]$ such that $n'+n''=n$, we have %\marginpar{Qua $n$ e' la lunghezza di tutto $b$, magari cambiare con $n'$ and $n''$ such that $n'+n''=n$}
  \begin{center}
    \tikzfig{tapes/cipriano/lemma1fperg}
  \end{center}
  \begin{center}
    \tikzfig{tapes/cipriano/lemma1fpergpt2}
  \end{center}
\end{proof}







\begin{lemma}\label{lemma:complete3}
  For all $b\in \Diag{\SigB}[1,A^n]$ and $c,d\in \Diag{PB}[A^n,A^m]$, 
  \begin{center} if $\osemPB{b;c}=\osemPB{b;d}$, then $\osemPB{b;D_c}=\osemPB{b;D_d}$ and $\osemPB{b;T_c}=\osemPB{b;T_d}$.\end{center}
\end{lemma}




\begin{theorem}\label{thm:completenesspartialcircuits}[Completeness]
  For all $c,d\in \Diag{PB}[A^n,A^m]$, if $\osemPB{c}=\osemPB{d}$ then $c=_{\mathbb{\SigPB }}d$.
\end{theorem}
\begin{proof}
  For every $b\in \Diag{\SigB}[1,A^n]$, 
  \begin{align}
    \osemPB{b;c}&=\osemPB{b};\osemPB{c} \notag\\
    &=\osemPB{b};\osemPB{d} \tag{Hp.}\\
    &=\osemPB{b;d}. \notag
  \end{align}
  Lemma~\ref{lemma:complete3} implies that, for every $b\in \Diag{\SigB}[1,A^n]$, $\osemPB{b;D_c}=\osemPB{b;D_d}$ and $\osemPB{b;T_c}=\osemPB{b;T_d}$, hence $\osemPB{D_c} = \osemPB{D_d}$ and $\osemPB{T_c} = \osemPB{T_d}$. Now, since $D_c,D_d,T_c$ and $T_d$ are Boolean circuits in $\Diag{\SigB}$, then by Theorem~\ref{thm:completenessbooleancircuits} it holds that $D_c\eqB D_d$ and $T_c\eqB T_d$. Finally, by Proposition~\ref{prop:decompositionpartialcircuits}, since $c$ and $d$ decompose through the same domain and total part, we have that $c\eqPB d$.
\end{proof}
Moreover, one can characterise exactly the image of $\Diag{\SigPB}$ through $\osemPB{-}$. Let $\Cat{Par}_2$ be the full subcategory of $\Cat{Par}$ where objects are powers of the set $2$, i.e.,  $2^n$ for all $n \in \mathbb{N}$.


\begin{proposition}\label{prop:isopartialboolean}
The monoidal functor  $\osemPB{-}\colon \Diag{{\SigPB }} \to \Cat{Par}$ factors as
\[\xymatrix{ \Diag{\SigPB } \ar@{->>}[r]& \Diag{\mathbb{\SigPB }} \ar[r]^{\cong}& \Cat{Par}_2 \ar@{^{(}->}[r]& \Cat{Par} }\]
where the leftmost and rightmost functors are the obvious quotients and injections and the central arrow is a monoidal isomorphism.
\end{proposition}



Theorem \ref{thm:completenesspartialcircuits} also provides a useful characterisation of diagrams in $\Diag{\SigPB}[1,A^n]$. First consider the diagram $\ncopierbis\colon A \to A^n$ defined inductively as
\begin{center}
  \tikzfig{tapes/cipriano/ncopierbis}
\end{center}
then define the diagram $\Flipn{\bot}\colon 1 \to n$ as $\Flip{\bot};\ncopierbis$.
\begin{lemma}\label{lemma:caratterizzazionecircuiti0n}
  For all $c\in \Diag{\SigPB}[1,A^n]$, either $c\eqPB \Flipn{\bot}$ or $c\eqPB b$ for some $b\in \Diag{\SigB}[1,A^n]$.
\end{lemma}
\begin{proof}
Recall that $\osemPB{c}$ is a partial function of type $1 \to 2^n$. Thus, it is either $\bot$ or  of the form $\osemPB{b}$ for some $b\in \Diag{\SigB}[1,A^n]$. In the first case, simple computations confirm that  $\osemPB{\Flipn{\bot}}=\bot$, thus $\osemPB{c}=\osemPB{\Flipn{\bot}}$. By Theorem~\ref{thm:completenesspartialcircuits}, $c\eqPB \Flipn{\bot}$. In the second case, $\osemPB{c}=\osemPB{b}$ and, again by Theorem~\ref{thm:completenesspartialcircuits}, $c\eqPB b$.
\end{proof}
%
%%% OLD
%
\begin{comment}
for every $y\in 2^n$, 
  \begin{align}
     \osem{\Flipn{\bot}}(y\vert \bullet)&=\osem{\Flip{\bot};\ncopierbis}(y\vert \bullet) \tag{Def.}\\
      &=(\osem{\Flip{\bot}}; \osem{\ncopierbis})(y\vert \bullet) \tag{Funct. of $\osem{-}$}\\
      &=\sum_{x\in 2} \osem{\Flip{\bot}}(x\vert \bullet)\cdot \osem{\ncopierbis}(y\vert x) \tag{; of $\KlD$}\\
        &=\sum_{x\in 2} \delta_{(x,x)}(1,0) \cdot \osem{\ncopierbis}(y\vert x) \tag{Def. of $\osem{\Flip{\bot}}$}\\
      &=0 \notag
      \end{align}
      hence, $\osem{f}=\osem{\Flipn{\bot}}$, and by Theorem~\ref{thm:completenesspartialcircuits} we have $f=\Flipn{\bot}$. In the second case, $\osem{f}=\osem{b}$ for some $b\in \Diag{\SigB}[1,A^n]$, then Theorem~\ref{thm:completenesspartialcircuits} implies $f=b$.
\end{comment}





%
%
%
% SUBSECTION 
%
%
\subsection{Probabilistic Boolean circuits}\label{ssec:probabilisticboolean}
Finally, we illustrate probabilistic Boolean circuits from \cite{piedeleu2025boolean}.
%
Consider the monoidal signature  \[ \SigPRB \defeq \SigPB \cup\{ \Flip{p} \mid p\in (0,1)\}\]  
 where $\Flip{p}$ denotes a probabilistic gate that receives no input and emits a Boolean signal on the right, which is $1$ with probability $p$ and $0$ with probability $1-p$. Thus, string diagrams in $\Diag{\SigPRB}$ represent arrows of $\KlD$. 


To formally define their semantics, recall the symmetric monoidal category $(\KlD, \otimes, \uno)$ and observe that there is an obvious monoidal embedding $\ParJ\colon \Cat{Par} \to \KlD$ which is the identity on objects and associates to a partial function $f\colon X \to Y$ the function $\ParJ(f)\colon X \to \Dis(Y)$, mapping $x$ to $\delta_{f(x)}$, if $f(x)$ is defined and to the null distribution $\star$ otherwise. %\marginpar{Se non si usa piu' $J$ cambierei la macro di $\ParJ$ con $J$ C: non mi sembra si usi più cambiamo la macro di ParJ e mettiamo dentro J?}




The semantics of $\SigPRB$ is given by the monoidal functor $\osem{-}\colon \Diag{\SigPRB} \to \KlD$ defined on objects as $A^n \mapsto 2^n$. On arrows, it is defined by first interpreting generators  $c\in \SigPB$ as $\ParJ(\osemPB{c})$, i.e.
\[
    \begin{array}{rclrclrcl}
      \osem{\Andgate}\colon 2 \times 2& \to & 2  & \osem{\Notgate}\colon  2& \to & 2 & \osem{\Flip{1}} \colon 1 & \to & 2 \\
      (x,y)&\mapsto&\delta_{x\wedge y} &  x & \mapsto & \delta_{\neg x} & \bullet &\mapsto&\delta_{1}\\[4pt]
       \osem{\CBcopier} \colon 2  & \to & 2\times 2 & \osem{\CBdischarger} \colon 2 & \to & 1 & \osem{\CBcocopier} \colon 2 \times 2 & \to & 2  \\
      x &\mapsto&\delta_{(x,x)} & x &\mapsto&\delta_{\bullet} & (x,y) &\mapsto& \begin{cases}\delta_{x} &\text{if }x=y\\ \star & \text{else}\end{cases}, 
    \end{array}
  \]

 the generator $\Flip{p}$ as 
\begin{equation}\label{eq:semanticsFlip}
\begin{array}{rcl}
      \osem{\Flip{p}}\colon 1& \to & 2   \\
      \bullet&\mapsto& \delta_1 +_p \delta_0  \end{array}
    \end{equation}
and then inductively on the structure of the monoidal category $(\KlD, \otimes, \uno)$. Simple computations confirm that $\osem{-}$ extends both $\osemB{-}$ and $\osemPB{-}$, in the sense that the following diagrams commute.
\[
\xymatrix{\Diag{\SigB} \ar@{^{(}->}[r] \ar[d]_{\osemB{-}}& \Diag{\SigPB} \ar@{^{(}->}[r] \ar[d]|{\osemPB{-}} & \Diag{\SigPRB} \ar[d]^{\osem{-}}\\
\Sets \ar@{^{(}->}[r]_{\Pa}& \Cat{Par} \ar@{^{(}->}[r]_{\ParJ}& \KlD
}
\]








\medskip
Figures~4 and~5 in \cite{piedeleu2025boolean} provide a sound and complete equational axiomatisation of probabilistic Boolean circuits. The system essentially consists of the axioms of Boolean algebras (reported in Figure~\ref{tab:booleanalgebra}), together with four non-trivial axioms governing the behaviour of $\Flip{p}$ and additional axioms for $\CBcocopier$, which differ from those presented in Figure~\ref{tab:partialbooleanalgebra}.

Importantly, the equivalence induced by these axioms is strictly coarser than the one given by $\osem{-}$. In fact, two diagrams $c$ and $d$ are equivalent under the axioms of \cite{piedeleu2025boolean} if and only if $\osem{c} \propto \osem{d}$. The relation $\propto$ is defined on arrows of $\KlD$ by setting $f \propto g$ whenever there exists a real number $\lambda > 0$ such that $f(y\mid x) = \lambda \cdot g(y\mid x)$ for all $x \in X$ and $y \in Y$. In the next section, we provide a complete  axiomatisation for the equivalence induced by $\osem{-}$ by means of probabilistic tape diagrams.



